\documentclass[10pt,a4paper]{amsart}
\usepackage[margin=19mm]{geometry}
\usepackage{amsmath,amssymb,mathtools}
\usepackage[scr=rsfs,cal=euler]{mathalfa}
\usepackage{xfrac}
\usepackage[unicode,hidelinks,bookmarks=false]{hyperref}
\newcommand{\ie}{i.e.\ }
\newcommand{\cf}{cf.\ }
\newcommand{\resp}{respectively\ }
\newcommand{\Subsection}{\subsection}
\providecommand{\nobreakdash}{\nobreak\hbox{-}}
\setlength{\emergencystretch}{2em}
\multlinegap0pt
\usepackage{amssymb}
\usepackage{mathtools}
\newtheorem{theorem}{Theorem}
\newcommand{\abs}[1]{\left|#1\right|}
\title{Permutations with a given $X$-descent set}
\author{Mohamed Omar}
\date{Source: arXiv:2402.10443v3 (CC BY 4.0)}
\begin{document}
\maketitle

\noindent\emph{Source and licence.} Permutations with a given $X$-descent set. Version arXiv:2402.10443v3 (CC BY 4.0), distributed by arXiv under the Creative Commons Attribution 4.0 International licence, http://creativecommons.org/licenses/by/4.0/, as recorded in the arXiv licence metadata for this exact version and shown on the abstract page https://arxiv.org/abs/2402.10443v3. Full licence notice: \texttt{licenses/arxiv-2402.10443-CC-BY-4.0.txt}.\par\smallskip

\noindent\textit{Editorial scope.} Counted unit: the article's theorem thm:recursive (source0.tex lines 222--236). The article's complete original proof of it is delivered with it (source0.tex lines 238--283, one \texttt{proof} environment of 46 lines). No mathematics is written, repaired or re-derived: the statement and the proof are the article's own lines, in the article's own notation, and no retained line is substituted or re-flowed.  Retained uncounted as the source-local context the proof needs: lines 169--171.  Nothing was omitted from any retained span, so the package carries no omission record.  No retained line contains a citation, so the wrapper carries no bibliography and no bibliography entry.  No retained line contains a figure environment, an image-inclusion command or drawing code, so no image is delivered and the package declares no asset dependency.  Editorial formatting only, in addition: an AMS \texttt{amsart} preamble that restates the article's own declarations and macros used by the retained lines, namely the environments \texttt{theorem} on separate counters, together with the standard packages those lines need.  The retained environments print in this document as Theorem 1 in order of appearance, whereas the article prints the same items with the numbering of its own full document.  The complete native source of the article as distributed in the arXiv e-print for this exact version is bundled unchanged as \texttt{sources/154660/source0.tex}.
\begin{equation}\label{eq:recursion-binomial}
d_X(I;n)=\binom{n}{m}\,d_X(I^-;m)\,d_X(\emptyset;n-m)-d_X(I^-;n).
\end{equation}

\begin{theorem}\label{thm:recursive}
Assume that $X$ is standardization-invariant.
Let $I=\{i_1<i_2<\cdots<i_k\}\subseteq[n-1]$, set $i_0=0$, and let $n\ge \max(I)+1$.
Then
\begin{equation}\label{eq:IE-expanded}
d_X(I;n)
=
\sum_{\{j_1<\cdots<j_r\}\subseteq\{1,\dots,k\}}
(-1)^{k-r}
\binom{n}{n-i_{j_r},i_{j_r}-i_{j_{r-1}},\ldots,i_{j_2}-i_{j_1}}
\prod_{t=0}^{r} d_X\!\left(\emptyset;\,i_{j_{t+1}}-i_{j_t}\right),
\end{equation}
where we interpret $j_0=0$ and $j_{r+1}=k+1$ with $i_{k+1}=n$.
(When $r=0$ the empty multinomial coefficient is $1$ and the last product is $d_X(\emptyset;n)$.)
\end{theorem}

\begin{proof}
We prove~\eqref{eq:IE-expanded} by induction on \(k=\abs{I}\), iterating the binomial recursion~\eqref{eq:recursion-binomial}.
If \(k=0\), then \(I=\emptyset\) and the sum in~\eqref{eq:IE-expanded} has only the empty choice \(r=0\), which gives
\(d_X(\emptyset;n)\).
If \(k=1\) and \(I=\{m\}\), the two subsets \(\emptyset\) and \(\{1\}\) recover~\eqref{eq:recursion-binomial} with
\(I^-=\emptyset\).

Now let \(k\ge 2\) and assume the formula holds for all sets of size \(k-1\).
We have \(I=\{i_1<i_2<\cdots<i_k\}\subseteq[n-1]\), with \(m=i_k\), and \(I^-=I\setminus\{m\}\).
Standardization invariance gives
\begin{equation}\label{eq:IE-induction-start}
 d_X(I;n)=\binom{n}{m}\,d_X(I^-;m)\,d_X(\emptyset;n-m)\ -\ d_X(I^-;n).
\end{equation}

Apply the induction hypothesis to \(d_X(I^-;m)\) and \(d_X(I^-;n)\).
For \(d_X(I^-;m)\), $I^-=\{i_1<i_2<\cdots<i_{k-1}\}$ and the ambient set has size $m$, so
\[
d_X(I^-;m)
=
\sum_{\substack{\{j_1<\cdots<j_r\}\subseteq [k-1]\\ 0\le r\le k-1}}
(-1)^{k-1-r}\,
\binom{m}{i_{j_r}}\binom{i_{j_r}}{i_{j_{r-1}}}\cdots\binom{i_{j_2}}{i_{j_1}}\,
\prod_{t=0}^{r} d_X(\emptyset;\,i_{j_{t+1}}-i_{j_t}),
\]
where we use the same conventions as in~\eqref{eq:IE-expanded} and interpret \(j_0=0\), \(i_{j_0}=0\), and
\(i_{j_{r+1}}=m\).
Similarly,
\[
d_X(I^-;n)
=
\sum_{\substack{\{j_1<\cdots<j_r\}\subseteq [k-1]\\ 0\le r\le k-1}}
(-1)^{k-1-r}\,
\binom{n}{n-i_{j_r},i_{j_r}-i_{j_{r-1}},\ldots,i_{j_2}-i_{j_1}}\,
\prod_{t=0}^{r} d_X(\emptyset;\,i_{j_{t+1}}-i_{j_t}),
\]
now with \(i_{j_{r+1}}=n\).

Substituting these two expansions into~\eqref{eq:IE-induction-start} produces two sums indexed by subsets of \([k-1]\).
In the first sum, the prefactor \(\binom{n}{m}\,d_X(\emptyset;n-m)\) simply appends \(m=i_k\) as the final selected
descent position and appends the final gap factor \(d_X(\emptyset;n-m)\).  Reindexing by adjoining \(k\) to the subset
\(\{j_1<\cdots<j_r\}\subseteq[k-1]\) therefore yields precisely the terms in~\eqref{eq:IE-expanded} indexed by subsets of
\([k]\) that contain \(k\), with the correct sign \((-1)^{k-(r+1)}=(-1)^{k-1-r}\).
The second sum contributes the terms indexed by subsets of \([k]\) that do not contain \(k\), and the extra minus sign
changes \((-1)^{k-1-r}\) into \((-1)^{k-r}\).
Together these two families of terms give exactly the full sum~\eqref{eq:IE-expanded}.
\end{proof}

\end{document}
